% Generated draft: requires mathematical and source review.
\documentclass[11pt]{article}
\usepackage[margin=25mm]{geometry}
\usepackage{fontspec}
\setmainfont{Latin Modern Roman}
\usepackage{amsmath,amssymb,mathtools,unicode-math}
\setmathfont{Latin Modern Math}
\usepackage{enumitem,needspace,xurl,hyperref}
\setlength{\parindent}{0pt}
\setlength{\parskip}{5pt}
\setlength{\emergencystretch}{3em}
\newcommand{\N}{\mathbb{N}}
\newcommand{\Z}{\mathbb{Z}}
\newcommand{\Q}{\mathbb{Q}}
\newcommand{\R}{\mathbb{R}}
\newcommand{\C}{\mathbb{C}}
\newcommand{\F}{\mathbb{F}}
\newcommand{\A}{\mathbb{A}}
\newcommand{\PP}{\mathbb{P}}
\newcommand{\E}{\mathbb{E}}
\newcommand{\eps}{\varepsilon}
\newcommand{\dd}{\,\mathrm{d}}
\begin{document}
\section*{UnsolvedMath ID: 30004373}
\textbf{Source problem code:} catalogue entry
\subsection{Definitions and mathematical statement}
The source does not specify the ground field or explicitly state the finite-dimensional algebra conventions. The following makes the question precise conditionally in the customary setting of finite-dimensional associative unital algebras over an algebraically closed field \(k\). These conventions require source confirmation; no characteristic restriction is supplied.

Let \(A\) be such an algebra. It is \emph{indecomposable} if it is not isomorphic to a direct product of two nonzero algebras. It is \emph{symmetric} if
\[
 A\cong D(A):=\operatorname{Hom}_k(A,k)
\]
as \(A\)-\(A\)-bimodules, where \((a\cdot f\cdot b)(x)=f(bxa)\).

For each integer \(d\geq 1\), let \(\mu_A(d)\) be the least number of \(A\)-\(k[t]\)-bimodules, each free of rank \(d\) as a right \(k[t]\)-module, whose specializations
\[
 M\otimes_{k[t]} k[t]/(t-\lambda),\qquad \lambda\in k,
\]
include, up to isomorphism, all but finitely many indecomposable left \(A\)-modules of \(k\)-dimension \(d\). Set \(\mu_A(d)=\infty\) if no such finite collection exists. The algebra is \emph{tame} if \(\mu_A(d)<\infty\) for every \(d\). It has \emph{polynomial growth} if there exist \(C>0\) and an integer \(r\geq0\) such that \(\mu_A(d)\leq C d^r\) for every \(d\geq1\). Thus, within the tame class, non-polynomial growth means
\[
 \forall C>0\ \forall r\in\Z_{\geq0}\ \exists d\geq1:
 \mu_A(d)>C d^r.
\]

For a finite-dimensional left \(A\)-module \(M\), let \(\Omega_A^n(M)\) denote the \(n\)-th syzygy in a minimal projective resolution. The module is \emph{periodic} if \(\Omega_A^n(M)\cong M\) for some integer \(n>0\). Put \(A^e=A\otimes_k A^{\mathrm{op}}\), so that \(A\)-\(A\)-bimodules are left \(A^e\)-modules. The algebra \(A\) is \emph{periodic} if \(\Omega_{A^e}^n(A)\cong A\) for some \(n>0\); its period is conventionally the least such \(n\).

\emph{Question, under the conditional conventions above.} For every indecomposable symmetric tame algebra \(A\) of non-polynomial growth, does
\[
 \bigl(\forall\text{ simple left }A\text{-modules }S\bigr)
 \bigl(\exists n_S\in\Z_{>0}\bigr)\quad
 \Omega_A^{n_S}(S)\cong S
\]
imply that \(A\) is periodic with least bimodule period \(4\)? Explicitly, the latter conclusion is
\[
 \Omega_{A^e}^{4}(A)\cong A,
 \qquad
 \Omega_{A^e}^{j}(A)\not\cong A\quad(1\leq j<4).
\]
The source says ``period \(4\)'' without defining whether minimality is intended. If it means only \(\Omega_{A^e}^{4}(A)\cong A\), the exclusion of smaller periods must be omitted; this distinction remains unresolved.
\subsection{Short English statement}
Must an indecomposable symmetric tame algebra of non-polynomial growth have bimodule period \(4\) whenever every simple module is periodic? The ground-field hypotheses and whether \(4\) must be the least period require confirmation from the source.
\subsection{Sources}
\begin{itemize}
\item Local UnsolvedMath registry, permanent ID \(30004373\), source problem code \texttt{catalogue entry}, ``Periodicity of Symmetric Tame Algebras,'' in the Oberwolfach Reports Open Problems collection. The supplied record assigns proposed year \(2020\) and names no individual proposer.
\item Supplied source citation: \emph{Representation Theory of Quivers and Finite Dimensional Algebras} (2021), Oberwolfach Reports, pp.\ 143--230, DOI \texttt{10.4171/owr/2020/3}. \url{https://doi.org/10.4171/owr/2020/3}. The supplied background describes Oberwolfach Reports as documenting mathematical workshops and their open-problem discussions.
\item Supplied literature source: \emph{Periodicity shadows I}. \url{https://arxiv.org/abs/2411.17381}.
\item Literature checked at: \texttt{2026-08-22}. This is supplied dated curation, not independently verified. Its status label is ``open,'' and its assessment says that period four is established in important families, remains conjectural in the nonpolynomial tame symmetric setting, and that current classification work still treats the general assertion as open. The supplied background requests independent verification of time-sensitive status and cited claims.
\item The registry labels the statement \texttt{corrected\_verified} and reports a manual semantic, typographic, and LaTeX review. These are supplied editorial assertions, not independent source verification or mathematical review of this draft.
\end{itemize}
\end{document}
